\documentclass[11pt]{article}

\usepackage[margin=1in]{geometry}
\usepackage{amsmath,amssymb,amsthm}
\usepackage{mathtools}
\usepackage{enumitem}
\usepackage[colorlinks=true,linkcolor=blue,citecolor=blue,urlcolor=blue]{hyperref}
\usepackage{xcolor}
\usepackage{tikz}
\usetikzlibrary{shapes,positioning}

% ---------------------------------------------------------------------
% Status-labelling discipline. Every mathematical claim below carries one
% of these five tags, defined here and used consistently throughout.
% ---------------------------------------------------------------------
\newtheorem{theorem}{Theorem}[section]
\newtheorem{proposition}[theorem]{Proposition}
\newtheorem{corollary}[theorem]{Corollary}
\newtheorem{lemma}[theorem]{Lemma}
\newtheorem{definition}[theorem]{Definition}
\newtheorem{extthm}[theorem]{External Theorem}
\newtheorem{conjecture}[theorem]{Conjecture}
\theoremstyle{remark}
\newtheorem{heuristic}[theorem]{Heuristic}
\newtheorem{falsified}[theorem]{Falsified}
\newtheorem{remark}[theorem]{Remark}
\newtheorem{example}[theorem]{Example}

\title{Where the Problem Stands: Contextual States on Concrete Orthomodular
Lattices\\[4pt]\large A hand-held mathematical exposition}
\author{}
\date{}

\begin{document}
\maketitle

\begin{center}
\begin{minipage}{0.92\textwidth}
\small
\noindent\textbf{How to read this document.} \textbf{Theorem / Proposition /
Corollary / Lemma}: proved in this programme's own work (proof given, marked
where compressed). \textbf{External Theorem}: proved in the literature, cited
not re-derived. \textbf{Conjecture}: precise, open, not asserted.
\textbf{Heuristic}: an unproved intuition. \textbf{Falsified}: investigated
and shown false, recorded so it is not retried under a new name.
\end{minipage}
\end{center}

\bigskip

\section*{Summary}
\addcontentsline{toc}{section}{Summary}

Can a concrete system of yes/no questions carry a probability assignment that
is consistent at every finite stage, yet has no countably-additive completion?
Formally: a two-valued state $s_0$ on a finite piece of a concrete
$\sigma$-class $L$ that is finitely coherent (agrees with \emph{some} global
finitely additive state) but extends to \emph{no} global $\sigma$-additive
one. If $L$ is only an \emph{orthomodular poset} (meets not required), the
answer is \textbf{yes} — Theorem~\ref{thm:main} below, ZFC, no extra
hypotheses. If $L$ must be an \emph{orthomodular lattice} (every pair has a
meet), the answer is \textbf{conjectured no} (Conjecture~\ref{conj:main}) and
open. Three obstruction results plus one external theorem all point toward
``no'': the obstacle is a genuinely infinitary, $\omega_1$-scale combinatorial
question (\S5), and no finite method can reach it. Two precise open routes are
stated in \S7.

\tableofcontents

% =======================================================================
\section{The setting, from scratch}
% =======================================================================

We build the objects the problem is about, slowly. Nothing here requires
quantum logic; only Boolean algebras, measures, and order theory. Each
definition gets a one-sentence gloss, a formal statement, and a computation.

\subsection{Why orthomodular lattices, informally}

Classical probability assumes every pair of events is compatible: the events
of an experiment form a Boolean algebra. Drop that assumption — keep only
``$A$ implies $B$'' and ``the complement of $A$,'' but stop requiring that any
two events have a joint refinement — and what survives is an
\emph{orthomodular lattice}. Everything below is pure mathematics about such
structures; no physics is assumed.

\subsection{Orthomodular lattices}

\begin{definition}[Orthomodular lattice]
\label{def:oml}
\textbf{Gloss:} events with a well-behaved ``did not happen'' operation, but
no distributive law.

\smallskip
\noindent A bounded lattice $L$ ($0,1,\wedge,\vee$) with an order-reversing
involution $a\mapsto a^\perp$ such that
\[
  a\vee a^\perp = 1,\qquad a\wedge a^\perp = 0,\qquad a\le b\Rightarrow
  b^\perp\le a^\perp,\qquad a^{\perp\perp}=a,
\]
and, the \emph{orthomodular law},
\[
  a\le b \ \Longrightarrow\ b = a\vee(b\wedge a^\perp).
\]
$L$ is Boolean exactly when every two elements are compatible (distributivity
holds); the interesting case is when they are not.
\end{definition}

\begin{example}[$MO_2$: the smallest non-Boolean case]
\label{ex:mo2}
Six elements: $0,1$, atoms $a,a^\perp,b,b^\perp$, with
\[
  a\wedge b = 0,\qquad a\vee b = 1,\qquad a\wedge b^\perp = 0,\qquad
  a\vee b^\perp = 1,
\]
symmetrically in $a^\perp,b^\perp$. Check orthomodularity directly: the only
nontrivial instance of $a\le b\Rightarrow b=a\vee(b\wedge a^\perp)$ has $a\le
b$ forcing $a\in\{0,b\}$, so it holds vacuously or trivially — no case
produces a violation. Distributivity fails at
\[
  a\wedge(b\vee b^\perp) = a\wedge 1 = a,\qquad\text{but}\qquad
  (a\wedge b)\vee(a\wedge b^\perp) = 0\vee 0 = 0,
\]
and $a\neq 0$. So $MO_2$ is orthomodular but not Boolean — the minimal such
example. Every finite orthomodular lattice, $MO_2$ included, turns out to be
\emph{tame} in the sense of Definition~\ref{def:tame} — Proposition
\ref{prop:finite-obstruction} below.
\end{example}

\begin{example}[The pentagon]
\label{ex:pentagon}
Five atoms $p_1,\dots,p_5$ in a cycle; $p_i$ compatible with $p_{i+1}$ (indices
mod 5, sharing a Boolean block $\{0,p_i,p_{i+1},p_i^\perp\wedge
p_{i+1}^\perp,1\}$), incompatible with $p_{i+2}$. Finite, concrete,
non-Boolean, and — used again in \S8 — a standard test bed precisely because
finiteness forces tameness regardless of what else is true of it.
\end{example}

\subsection{Concrete (set-representable) orthomodular lattices}

\begin{definition}[Concrete $\sigma$-class; OMP vs.\ OML]
\label{def:concrete}
\textbf{Gloss:} subsets of a set $\Omega$, closed under complement and
countable disjoint union; a lattice exactly when every pair also has a
genuine set-theoretic meet inside the family.

\smallskip
\noindent Fix $\Omega$. A \emph{concrete $\sigma$-class} $L\subseteq
\mathcal P(\Omega)$ contains $\emptyset,\Omega$, is closed under
complement $E\mapsto\Omega\setminus E$, and closed under countable unions of
pairwise-disjoint members. Order is $\subseteq$. Every such $L$ is
automatically an \textbf{OMP} (orthomodular poset — pairs need not have a
meet in $L$). If additionally every $A,B\in L$ have a $\subseteq$-greatest
lower bound in $L$, call $L$ an \textbf{OML} (orthomodular lattice). The OMP
vs.\ OML distinction is the entire subject of this document.
\end{definition}

Concreteness matters: an abstract orthomodular lattice can be represented in
exotic non-set ways, but a concrete one is tied to honest membership facts.
(The projection lattice of a Hilbert space generically fails to be concrete —
Kochen--Specker — but every object below is concrete by definition, so that
obstruction never enters here.)

\subsection{States}

\begin{definition}[Two-valued state]
\label{def:state}
\textbf{Gloss:} a yes/no probability, additive over disjoint unions.

\smallskip
\noindent For $B\subseteq L$ closed under complement, a state is
$s:B\to\{0,1\}$, $s(\Omega)=1$, with
\[
  s\Big(\bigcup_i A_i\Big) = \sum_i s(A_i)
\]
for every pairwise-disjoint $\{A_i\}\subseteq B$ with union in $B$: finite
family for \emph{finitely additive}, countable family for
\emph{$\sigma$-additive}. For $\omega\in\Omega$, $\delta_\omega(E) :=
\mathbf 1[\omega\in E]$ is always $\sigma$-additive — a \emph{Dirac} state;
every other state is \emph{non-Dirac}.
\end{definition}

The entire problem lives in the gap between ``additive over finite unions''
and ``additive over countable unions.''

\subsection{Blocks}

\begin{definition}[Block]
\label{def:block}
\textbf{Gloss:} the largest all-compatible piece containing a given event —
an ordinary Boolean $\sigma$-algebra sitting inside $L$.

\smallskip
\noindent A \emph{block} of $L$ is a maximal pairwise-compatible subfamily.
Every element of $L$ lies in some block (Zorn's lemma), and each block, on
its own, is a Boolean $\sigma$-algebra of subsets of $\Omega$.
\end{definition}

All non-classicality lives \emph{between} blocks, never inside one — the
single fact the rest of this document is built on (made precise as Theorem
\ref{thm:blocks-boolean}).

\subsection{Essential contextual states}

\begin{definition}[Essential contextual state]
\label{def:essential}
\textbf{Gloss:} a finite probability fragment consistent at the finite level,
yet incompatible with \emph{every} countably-additive completion.

\smallskip
\noindent Fix a concrete $\sigma$-class $L$, a finite $\perp$-closed
$B\subseteq L$, and a state $s_0$ on $B$. Call $s_0$ \emph{finitely coherent
on $L$} if it extends to some finitely additive state on all of $L$. Call
$s_0$ \emph{$\sigma$-essential} if it is finitely coherent on $L$ but no
$\sigma$-additive state on $L$ extends it.
\end{definition}

The finite-coherence clause is load-bearing, not decorative — quantified
precisely in Proposition~\ref{prop:amendment}: without it, a
self-contradictory fragment that no state of any kind extends would trivially
(and vacuously) satisfy the definition.

\begin{definition}[Tame]
\label{def:tame}
$L$ is \emph{tame} ($\Phi(L)$) if every finitely coherent finite pattern on
$L$ extends to a $\sigma$-additive state — i.e.\ $L$ carries no
$\sigma$-essential state.
\end{definition}

% =======================================================================
\section{The central question, stated}
% =======================================================================

\begin{conjecture}[The central conjecture]
\label{conj:main}
Every concrete $\sigma$-complete orthomodular \textbf{lattice} is tame.
\end{conjecture}

\begin{center}
\begin{tabular}{ll}
\hline
OMP (Definition~\ref{def:concrete}, no meet requirement) & \textbf{false} — witness exists, Thm.~\ref{thm:main}\\
OML (meets required) & \textbf{open} — Conjecture~\ref{conj:main}\\
\hline
\end{tabular}
\end{center}

\subsection{A second, equivalent, face}

\begin{theorem}[Cluster reformulation]
\label{thm:phi-cluster}
Call $\mathcal C\subseteq L$ finite an \emph{f.a.-coherent cluster} if some
finitely additive state on $L$ assigns $1$ to every member. Then:
\begin{enumerate}[label=(\alph*)]
\item (block form) every finite pattern on a block, coherent with some
  finitely additive state on $L$, extends to a $\sigma$-additive state on
  $L$;
\item (cluster form) every f.a.-coherent cluster $\mathcal C$ is assigned $1$
  throughout by a single $\sigma$-additive state.
\end{enumerate}
are equivalent.
\end{theorem}

\emph{Proof.} (a)$\Rightarrow$(b): given $\mathcal C$, let $\mu$ witness its
coherence; the pattern $s_0 := \mu|_{\perp\text{-closure}(\mathcal C)}$ is a
block-level finite pattern coherent with $\mu$, so (a) extends it
$\sigma$-additively, and that extension assigns $1$ throughout $\mathcal C$
by construction. (b)$\Rightarrow$(a): given a coherent block pattern $s_0$ on
$B$, let $\mathcal C := \{A\in B: s_0(A)=1\}$; this is an f.a.-coherent
cluster, so (b) gives a $\sigma$-additive state assigning it $1$ throughout,
and $\perp$-closure forces the same state to assign the $s_0$-false members
$0$, so it agrees with $s_0$ on all of $B$. $\hfill\square$

Both directions use no property of $L$ beyond $\perp$-closure — in
particular, no latticehood. Machine-verified in both directions.

We do \emph{not} present a third independent face. The tempting split ``no
Dirac state realizes the pattern'' / ``no non-Dirac state does either'' is a
tautology (every $\sigma$-state is one or the other), not a reformulation —
its real content is Theorem~\ref{thm:localization} in \S3.

% =======================================================================
\section{What is proved: the solid ground}
% =======================================================================

\subsection{All non-classicality lives between blocks}

\begin{theorem}[Blocks are Boolean $\sigma$-fields]
\label{thm:blocks-boolean}
Every block $M$ of a concrete $\sigma$-class $L$, restricted to itself, is a
Boolean $\sigma$-algebra of subsets of $\Omega$: closed under complement and
countable disjoint union, with every such union already inside $M$.
Consequently a countable disjoint family is tested for $\sigma$-additivity
inside the single block containing it.
\end{theorem}

\emph{Proof sketch.} Maximality of $M$ as a pairwise-compatible family forces
closure under the Boolean operations restricted to $M$ (standard
maximality argument). $\hfill\square$

\emph{What this buys.} The search for a witness reduces entirely to an
\emph{assembly} question — how finite patterns on different blocks can or
cannot be jointly realized — never a within-block question.

\subsection{Why finite coherence is load-bearing}

\begin{proposition}[Boolean baseline]
\label{prop:amendment}
If $L$ is (all of) a Boolean $\sigma$-algebra and $s_0$ is finitely coherent
on $L$, then $s_0$ extends to a Dirac state on $L$. So a Boolean carrier is
tame for free, not merely by conjecture.
\end{proposition}

\emph{Proof.} Let $\mu\supseteq s_0$ witness coherence, $B_0$ the finite
Boolean subalgebra generated by $B$. Its atoms partition $\Omega$ into
finitely many pieces $a_1,\dots,a_k$; finite additivity gives
\[
  1 = \mu(\Omega) = \sum_{j=1}^k \mu(a_j),
\]
each $\mu(a_j)\in\{0,1\}$, so exactly one $a_{j_0}$ has $\mu(a_{j_0})=1$ and
$a_{j_0}\neq\emptyset$. Any $\omega\in a_{j_0}$ gives $\delta_\omega$ agreeing
with $\mu$ on $B_0$, since each member of $B_0$ either contains $a_{j_0}$ or
is disjoint from it. $\hfill\square$

\emph{Why it matters.} Drop finite coherence and a self-contradictory pattern
(no global finitely additive state agrees with it at all) trivially fails to
extend $\sigma$-additively — for a reason with nothing to do with countable
additivity. Proposition~\ref{prop:amendment} shows the clause is not merely
excluding degenerate cases: on the \emph{most} classical carrier, finite
coherence already forces the strongest possible outcome. Any genuine witness
must live where this forcing fails — non-Boolean, as \S4's theorem does.

\subsection{Countably generated blocks force triviality}

\begin{theorem}[Countably generated boundary extension]
\label{thm:ctbly-gen}
Let $A\subseteq L$ be a Boolean $\sigma$-subalgebra generated by countably
many $(a_n)_{n<\omega}$, faithfully $\sigma$-embedded into a
$\sigma$-complete OML $L$ whose $\sigma$-states order-separate points. Then
every $\sigma$-state $\mu$ on $A$ extends to one on $L$.
\end{theorem}

\emph{Proof.} Set $b_n := a_n$ if $\mu(a_n)=1$, else $b_n:=a_n^\perp$; each
finite meet $b_1\wedge\cdots\wedge b_n$ has $\mu$-value $1$. $\sigma$-additivity
of $\mu$ forces
\[
  c_\mu := \bigwedge_{n<\omega} b_n \ \neq\ 0.
\]
Order-separation in $L$ supplies a $\sigma$-state on $L$ charging the image
of $c_\mu$; it agrees with $\mu$ on each generator $a_n$, hence on all of
$A$. $\hfill\square$

\emph{What this buys.} Any witness confined to — or faithfully,
separatingly embedded in — a countably generated piece always extends. A
witness, if one exists on a lattice, cannot be countably generated in this
sense.

\subsection{The Dirac/non-Dirac partition}

\begin{theorem}[Localization to the non-Dirac core]
\label{thm:localization}
For finite $\perp$-closed $B\subseteq L$ and state $s_0$ on $B$: $s_0$
extends to no $\sigma$-state on $L$ $\iff$
\[
  \underbrace{\bigcap\{A\in B : s_0(A)=1\} = \emptyset}_{\text{(i) no Dirac
  extension}} \quad\text{and}\quad \underbrace{\text{no non-Dirac
  $\sigma$-state extends $s_0$.}}_{\text{(ii)}}
\]
\end{theorem}

\emph{Proof.} Every $\sigma$-state is Dirac or non-Dirac, mutually exclusive
and jointly exhaustive, so ``none extends'' $\iff$ (i) and (ii). For (i):
$\delta_\omega|_B = s_0$ iff $\omega$ lies in every $s_0$-true member of $B$
(the $s_0$-false members are automatic by $\perp$-closure). $\hfill\square$

This is a tautology — ``a state is Dirac or it isn't'' — not an independent
reformulation. Its use: (i) is a purely finite, checkable condition that
Theorem~\ref{thm:main}'s construction shows can be arranged \emph{independently}
of (ii). All difficulty collapses onto (ii).

\subsection{The necessary condition, with its exact choice-strength}

\begin{theorem}[Non-extendability $\Leftrightarrow$ FIP, calibrated]
\label{thm:nonextend}
For $s_0$ on finite $\perp$-closed $B$, let $V := \{A\in B : s_0(A)=1\}$.
TFAE:
\begin{enumerate}[label=(\arabic*)]
\item some $\omega\in\bigcap V$ (Dirac realization);
\item $s_0$ extends to the Boolean $\sigma$-algebra generated by $B$;
\item $V$ has the \emph{finite intersection property} (FIP): every finite
  subfamily of $V$ has nonempty intersection.
\end{enumerate}
$(1)\Rightarrow(3)$ and $(2)\Rightarrow(3)$: plain ZF. $(3)\Rightarrow(2)$:
exactly the Boolean Prime Ideal theorem (BPI) — every proper filter extends
to an ultrafilter — and no weaker principle suffices.
\end{theorem}

\emph{Proof.} $(1)\Rightarrow(3)$: $\omega$ witnesses every finite
subfamily's intersection. $(2)\Rightarrow(3)$: an extension is finitely
additive and two-valued on the generated algebra, so finite meets of
value-$1$ sets stay value-$1$, hence nonempty. $(3)\Rightarrow(2)$: FIP makes
$V$ a filter base; BPI extends it to an ultrafilter on the generated algebra;
read off the two-valued state. $\hfill\square$

\begin{corollary}
\label{cor:stage3}
If $s_0$ is $\sigma$-essential, then (no Dirac realizes it, by definition)
$V$ fails FIP: some finite subfamily of $V$ has empty intersection. This
consequence is a plain ZF fact — the contrapositive of $(1)\Leftrightarrow(3)$
needs no choice; only the abstract three-way equivalence sits at BPI.
\end{corollary}

\emph{What this buys.} Any witness search reduces to finite value-$1$
families failing FIP at some checkable finite stage. Checked directly against
Theorem~\ref{thm:main}: its value-$1$ family fails FIP at
\[
  |V_{\text{crit}}| = 3,
\]
the smallest possible stage. Reverse-math reading: the \emph{obstruction}
(finite-stage FIP failure, used as a necessary condition) is choice-free;
only the \emph{positive} extension direction needs BPI and nothing stronger.

\subsection{The finite obstruction results}

\begin{proposition}[No finite witness; the two-cell obstruction]
\label{prop:finite-obstruction}
No finite concrete OMP carries a $\sigma$-essential state (finite $\Rightarrow$
every finitely additive state is automatically $\sigma$-additive — nothing
infinite to fail to converge on). Further: no concrete $\sigma$-complete
\textbf{OML} can contain two Boolean blocks that (a) each faithfully,
separatingly represent the same decreasing clopen sequence $(U_n)_{n<\omega}$
converging to a point, and (b) send that sequence to \emph{different} points
$i\neq j$ in the two blocks — because
\[
  \bigwedge_{n<\omega} U_n
\]
is computed once, in the ambient lattice, and cannot equal $\{i\}$ as seen
from one block and (something containing) $\{j\}\neq\{i\}$ as seen from the
other.
\end{proposition}

\emph{Proof.} The countable meet of a fixed decreasing family is an intrinsic
feature of a $\sigma$-complete lattice, computed via complementation and
countable-join closure; both faithful embeddings realize the \emph{same}
ambient meet, so it cannot disagree with itself. $\hfill\square$

\emph{What this buys.} Rules out the most natural gluing construction — two
Cantor-space-like blocks glued along a shared countable boundary — regardless
of where the gluing sits in the structure.

\begin{proposition}[The countably-generated route is fully closed]
\label{prop:ctbly-gen-full}
Combine Theorem~\ref{thm:ctbly-gen} with the Lusin--Souslin injection theorem
(countably many countably-generated quotient algebras of a standard Borel
space that jointly separate points already generate the full Borel
$\sigma$-algebra): no countably-generated, faithfully separating boundary
structure — however many countably-generated pieces are combined — supports
a witness.
\end{proposition}

\emph{What this buys.} Any surviving construction must be genuinely,
unavoidably uncountably generated.

% =======================================================================
\section{The one existing witness: a proved ZFC theorem}
% =======================================================================

\begin{theorem}[A $\sigma$-essential state exists, in ZFC]
\label{thm:main}
There is a concrete $\sigma$-complete orthomodular \textbf{poset} $L$ on
$\Omega$, $|\Omega|=\aleph_1$, carrying a $\sigma$-essential state. Plain
ZFC; no large cardinal or independence hypothesis.
\end{theorem}

\emph{Construction.} Let $M$ have $|M|=\aleph_1$, well-ordered with every
initial segment countable.

\smallskip
\noindent\textbf{Step 1 (Ulam matrix).} Build $\{C_{\alpha,n} :
\alpha<\omega_1,\ n<\omega\}\subseteq \mathcal P(M)$ with
\[
  \bigsqcup_{n<\omega} C_{\alpha,n} = M \setminus (\text{countable set}),
  \qquad
  C_{\alpha,n}\cap C_{\beta,n} = \emptyset\ (\alpha\neq\beta).
\]
(Each row disjoint and co-countable; each column pairwise disjoint across
rows.)

\smallskip
\noindent\textbf{Step 2 (rigidity).} Let $\Omega = M\times\{1,2,3,4\}$,
$D_{\alpha,n} := C_{\alpha,n}\times\{1,2,3,4\}$, and let $L$ be generated by
all singletons, the cells $D_{\alpha,n}$, and the three sets of Step 3. Claim:
every non-Dirac candidate fails. If a $\sigma$-state kills every singleton
(hence, by $\sigma$-additivity, every countable set), then in each row
exactly one cell has value $1$. The map
\[
  \alpha \ \longmapsto\ n(\alpha) := \text{the unique } n \text{ with }
  s(D_{\alpha,n})=1
\]
sends $\omega_1\to\omega$; since $|\omega_1|>\aleph_0$, some fiber is
uncountable, so pick $\alpha\neq\beta$ with $n(\alpha)=n(\beta)=n^\ast$. Then
$D_{\alpha,n^\ast}\cap D_{\beta,n^\ast}=\emptyset$ (column-disjointness), so
additivity forces
\[
  s\big(D_{\alpha,n^\ast}\sqcup D_{\beta,n^\ast}\big) = s(D_{\alpha,n^\ast}) +
  s(D_{\beta,n^\ast}) = 2,
\]
impossible for a $\{0,1\}$-valued state. Hence no non-Dirac $\sigma$-state
exists: every $\sigma$-state on $L$ is Dirac.

\smallskip
\noindent\textbf{Step 3 (the pattern).} Let
\[
  A := M\times\{1,2\},\qquad B := M\times\{1,3\},\qquad C := M\times\{2,3\}.
\]
Then
\[
  A\cap B = M\times\{1\},\qquad A\cap C = M\times\{2\},\qquad B\cap C =
  M\times\{3\},
\]
all nonempty, but
\[
  A\cap B\cap C = \emptyset.
\]
So $V=\{A,B,C\}$ fails FIP at stage exactly $3$ — matching Corollary
\ref{cor:stage3} at the minimum possible value. Fix a non-principal
ultrafilter $\mathcal U$ on $M$ and define $m$ on $L$ by voting through
$\mathcal U$; a short parity computation on the four fibers shows $m(A)=m(B)=
m(C)=1$ simultaneously, consistently, with no point witnessing all three.
This $m$ is a global finitely additive state, so $s_0 := m|_{\perp
\text{-closure}(A,B,C)}$ is finitely coherent.

\smallskip
\noindent\textbf{Conclusion.} $s_0$ is finitely coherent (by $m$); by
rigidity (Step 2) the only candidate $\sigma$-extensions are Dirac; no
$\omega\in\Omega$ lies in $A\cap B\cap C=\emptyset$. So no $\sigma$-state
extends $s_0$: it is $\sigma$-essential. $\hfill\square$

\begin{proposition}[Not a lattice]
\label{prop:not-lattice}
In $L$, $A\wedge B$ does not exist: every member of $L$ inside $A\cap B =
M\times\{1\}$ is countable (only countable sets and their complements/finite
combinations lie in $L$'s generated structure at that fiber), and among the
countable subsets of the uncountable set $M\times\{1\}$ there is no largest
one — given any countable $S\subsetneq M\times\{1\}$, adjoin one more point
of $M\times\{1\}\setminus S$ to get a strictly larger member of $L$.
\end{proposition}

\emph{Why this is the crux.} Theorem~\ref{thm:main} resolves the poset form
completely and says nothing about the lattice form, because its carrier is
not a lattice. Every later section is about this one gap.

% =======================================================================
\section{The external boundary result: measure extension on operator
algebras}
% =======================================================================

\begin{extthm}[Escolano--Peralta--Villena, 2025]
\label{thm:epv}
Let $J$ be a JBW$^\ast$-algebra with projection lattice $P(J)$, and suppose
$J$ has no direct summand of type $I_2$ (the spin-factor case, the
operator-algebraic analogue of $MO_2$). Then every bounded finitely additive
measure on $P(J)$ extends to a bounded linear functional on $J$ (and, via
Hahn--Banach, every Banach-space-valued one extends to a bounded operator).
\emph{Sharp:} if $J$ has a type-$I_2$ summand, there is a bounded finitely
additive measure on $P(J)$ that does \emph{not} extend.
\end{extthm}

\emph{Mechanism, in one line.} Away from type $I_2$, every bounded finitely
additive measure on $P(J)$ is automatically uniformly continuous — via
projection halving, small symmetric motions, and isoclinic interpolation —
and uniform continuity upgrades to linear extendability by standard
functional analysis. This is a \emph{finite}, elementary mechanism: no
uncountable cardinals, ultrafilters, or countable additivity in the sense of
Definition~\ref{def:state} enter.

\begin{remark}[What it does and does not tell us]
\label{rem:epv-scope}
\emph{Does:} narrows where an analogous phenomenon could live — the
symmetry-poor, type-$I_2$-rich regime, nowhere else.

\emph{Does not:} (1) it is a statement about JBW$^\ast$ projection lattices,
which need carry no compatible norm or Jordan structure — the concrete
$\sigma$-classes of Definition~\ref{def:concrete} are not, in general, of
this type, so the theorem neither forbids nor supplies a witness on them
directly; (2) the sharp (dual) direction's own mechanism is a finite
rank-counting clash inside a $3$-dimensional spin factor —
\[
  \tfrac13 \ne \tfrac12
\]
— categorically not the $\omega_1$-scale infinitary mechanism
Conjecture~\ref{conj:main} would need if false; (3) it narrows a
\emph{location}, not a proof — it supplies no argument for whether a witness
exists there.
\end{remark}

In one sentence: Theorem~\ref{thm:epv} narrows a target in an adjacent world;
it is not a step toward a proof here.

% =======================================================================
\section{Where the difficulty actually is: the honest crux}
% =======================================================================

A lattice-form witness, if one exists, must satisfy all four at once:
\[
\begin{array}{ll}
\text{(i)} & \text{uncountably generated (Thm.~\ref{thm:ctbly-gen}, Prop.~\ref{prop:ctbly-gen-full})}\\
\text{(ii)} & \text{symmetry-poor (Thm.~\ref{thm:epv}, heuristic pointer only)}\\
\text{(iii)} & \text{FIP-failing at some finite stage, yet finitely coherent (Thm.~\ref{thm:nonextend})}\\
\text{(iv)} & \text{survives lattice closure}
\end{array}
\]

\subsection{Why (iv) is the genuine obstacle}

Theorem~\ref{thm:main} achieves (i)--(iii) with room to spare and provably
fails (iv) (Proposition~\ref{prop:not-lattice}). Requiring $L$ to be a lattice
means: for \emph{every} pair, a meet must be manufactured and adjoined
consistently with the orthomodular law and all prior additivity constraints.
Doing this coherently for the uncountably many pairs forced by (i) is an
$\omega_1$-level combinatorial question: can a coherent system of forced
meets and joins be built across an uncountable index set without collapsing
the essential pattern to a single point?

\begin{remark}[Why finite methods cannot reach this]
\label{rem:finite-methods}
Every finite construction tried has produced only tame objects — forced, not
unlucky, by Proposition~\ref{prop:finite-obstruction}: finiteness makes
$\sigma$-additivity automatic, so no finite object can ever exhibit (iv)
failing at the uncountable scale (i) demands.
\end{remark}

\subsection{The current assessment}

\begin{remark}[Where the evidence leans — assessment, not theorem]
\label{rem:assessment}
Evidence leans toward Conjecture~\ref{conj:main} being \textbf{true}:
\begin{enumerate}[label=(\alph*)]
\item the two-cell countable-meet obstruction (Prop.~\ref{prop:finite-obstruction});
\item the countably-generated closure (Thm.~\ref{thm:ctbly-gen} + Prop.~\ref{prop:ctbly-gen-full});
\item every attempted uncountable ``hub'' construction has, on inspection,
  either collapsed back to the classical center or destroyed
  order-separation;
\item every candidate lattice-theoretic survival mechanism has, on
  inspection, either needed exactly the distributivity being ruled out, or
  been falsified outright (\S8).
\end{enumerate}
If proved, Conjecture~\ref{conj:main} becomes: latticehood alone forces every
finitely coherent pattern to extend $\sigma$-additively — a clean structural
theorem independent of any construction above.
\end{remark}

\begin{remark}[Which way effort is now pointed — strategic assessment, not a
theorem]
\label{rem:inversion}
A note on \emph{research allocation}, recorded so a reader knows how the two
routes of \S6 are being worked. Until recently the effort was organized as a
search for a witness, with ``prove no witness exists'' held as the fallback.
As of 2026-07-20 that ordering is inverted: \textbf{proving
Conjecture~\ref{conj:main} is now the primary target}, and the search for a
witness is retained as the standing falsification test. The reason is
bookkeeping, not a new theorem: three campaigns of attempted constructions
have each produced an obstruction rather than an object, so the accumulated
yield sits on the obstruction side. Nothing in \S3 or \S6 changes; the two
open routes are unchanged, and either one's resolution in either direction is
still genuine progress. What changes is which end is attacked first. This
remark asserts no mathematics: the conjecture remains open, and the
possibility that a witness exists remains entirely live --- were the evidence
strong enough to settle it, it would be stated above as a theorem and not
here as an assessment.
\end{remark}

% =======================================================================
\section{The two open directions, precisely}
% =======================================================================

\subsection{Route (a): does the countable-meet obstruction generalize?}

\begin{conjecture}[Generalized countable-meet obstruction]
\label{conj:route-a}
Does Proposition~\ref{prop:finite-obstruction}'s two-cell mechanism forbid
\emph{every} faithful, order-separating, countably-generated hub shared
between two blocks of a concrete $\sigma$-complete OML — not just the
specific Cantor-space gluing?
\end{conjecture}

If proved: eliminates every countably-generated-boundary construction at
once. If it fails: the surviving instance is a concrete design template for
a witness attempt.

\subsection{Route (b): the distributed nonseparating quotient route}

\begin{conjecture}[Distributed nonseparating quotient route]
\label{conj:route-b}
Is there a concrete $\sigma$-complete OML, built from uncountably-generated
boundary data whose $\sigma$-states do \emph{not} order-separate points
across the whole structure — evading Theorem~\ref{thm:ctbly-gen}'s
hypothesis by distributing nonseparating quotient copies so no global
$\sigma$-state assembles, while every finite and countable sub-piece stays
consistent?
\end{conjecture}

If constructible as a genuine lattice: refutes Conjecture~\ref{conj:main}
outright. If provably impossible: the complementary half of a full proof,
alongside (a).

Read together: (a) attacks the countably-generated case in full generality;
(b) attacks the uncountably-generated nonseparating case (a) does not cover.
A witness, if one exists, must refute (b) constructively while already
escaping (a)'s mechanism.

% =======================================================================
\section{What is heuristic, not yet mathematics}
% =======================================================================

\begin{falsified}[Dimensional-equivalence as a proxy for shared-boundary
agreement]
\label{fals:perspectivity}
Tested hypothesis: the lattice-theoretic analogue of Murray--von Neumann
equivalence of projections — call two atoms $p,q$ \emph{equivalent} if
related by a dimension-preserving symmetry of the lattice — might proxy for
``forced to agree under every state assigning both value $1$ on a shared
boundary'' (the relation actually doing work in \S3). \textbf{Falsified}: on
the pentagon (Example~\ref{ex:pentagon}) and every finite homogeneous OML,
\[
  p \sim q \quad \text{for \emph{every} pair } p,q
\]
— a purely structural consequence of finite homogeneity (any two atoms are
joined by a chain of complementary pairs). A relation identically true on
every pair cannot proxy a relation that must be selective. The two relations
are also of different type: one rank-sensitive, one rank-blind. Dead; not to
recur under a new name.
\end{falsified}

\begin{falsified}[A winding/parity invariant]
\label{fals:winding}
Built directly on Claim~\ref{fals:perspectivity}: a hoped-for integer
invariant tracking how forced meets/joins wind around a cycle of incompatible
atoms, meant to detect whether requirement (iv) could hold. Falls with its
foundation — an identically-true relation carries no invariant. Independently:
the two available test objects are structurally incapable of exhibiting the
phenomenon at all — the pentagon is finite, hence automatically tame
(Prop.~\ref{prop:finite-obstruction}); Theorem~\ref{thm:main}'s witness is
provably \emph{not} a lattice (Prop.~\ref{prop:not-lattice}), hence has no
forced meets for a winding invariant to be about. Two independent reasons the
idea dies.
\end{falsified}

\begin{heuristic}[A legitimate, unexplored forwarding lead]
\label{heur:z2states}
$\mathbb Z_2$-valued states (valued in the two-element group, not $\{0,1\}$
read as reals) on symmetric-difference-closed orthomodular structures are a
real, previously studied class in the literature — distinct in kind from the
falsified ad hoc invariant of Claim~\ref{fals:winding}, which does not appear
in the literature under any name. Whether they bear on the essential-witness
question is entirely open and unexplored: named here as a place to look
next, not as a result or even a formulated conjecture.
\end{heuristic}

% =======================================================================
\section*{Closing remark}
\addcontentsline{toc}{section}{Closing remark}

A witness to genuinely contextual, $\sigma$-inextendable probability data is
constructed and proved, unconditionally, in ZFC (Theorem~\ref{thm:main}) —
one order-theoretic notch below the structure the field cares about (poset,
not lattice). Whether latticehood alone forces every such pattern to dissolve
is Conjecture~\ref{conj:main}, open. Three obstruction results and one
external theorem in an adjacent world point toward ``yes, it dissolves,'' but
the obstruction, if it exists, is a genuinely infinitary $\omega_1$-question
(\S5) that no finite method can reach even in principle. Two precise routes
(\S6) remain: either one's resolution, in either direction, is genuine
progress.

\end{document}
